%% ma: L12-B07-0024
%% lop: 12
%% chuong: 2
%% bai: 7
%% dang: 12.7.9
%% loai: TN4
%% muc: NB
%% dap_an: "B"
%% nguon: "CĐ2.2.pdf giáo viên gửi 27/10/2026 (M07, Lê Minh Khôi – Nguyễn Anh Khôi, Chương II Chủ đề 2 Tọa độ của vectơ trong không gian), Câu 21"
%% kien_thuc: "Bài 7 (tọa độ điểm), Bài 8 (trung điểm)"
%% trang_thai: chinh-thuc
%% ghi_chu: "Giáo viên đã duyệt 27/10/2026"
\begin{ex}
Trong không gian $Oxyz$, cho điểm $X\left(-\dfrac32;1;-1\right)$. Gọi $X_1$ là điểm đối xứng của $X$ qua điểm $O$. Tọa độ của điểm $X_1$ là
\choice{$X_1\left(\dfrac32;1;1\right)$}{\True $X_1\left(\dfrac32;-1;1\right)$}{$X_1\left(-\dfrac32;1;1\right)$}{$X_1\left(-\dfrac32;-1;-1\right)$}
\loigiai{Đối xứng qua gốc $O$: đổi dấu cả ba tọa độ: $X_1\left(\dfrac32;-1;1\right)$.}
\end{ex}
